% الجزء: sets-functions-relations
% الفصل: relations
% القسم: relations-as-sets

\documentclass[../../../include/open-logic-section]{subfiles}

\begin{document}

\olfileid[ar]{sfr}{rel}{set}
\olsection{العلاقات بوصفها مجموعات}

\begin{explain}
من المجموعات المهمة المذكورة في \olref[sfr][set][imp]{sec}
$\Nat$ و$\Int$ و$\Rat$ و$\Real$. وبين !!{element}s هذه المجموعات
علاقات مألوفة؛ فلكل منها \emph{علاقة ترتيب} معروفة، وثمة علاقة
\emph{«مطابق لـ»} بين كل شيء ونفسه دون غيره. وسنعرض علاقات كثيرة
جديرة بالنظر، والممكن منها أكثر مما سنعرض. فلنبدأ، قبل تفصيلها،
ببيان أن العلاقات يمكن أخذها ضربًا مخصوصًا من المجموعات.

ويحتاج ذلك إلى أمرين من \olref[sfr][set][pai]{sec}. أحدهما
\emph{الزوج المرتب}: فمن $\OLId{latin-ordinary}{a}$ و$\OLId{latin-ordinary}{b}$ نصنع~$\tuple{\OLId{latin-ordinary}{a}, \OLId{latin-ordinary}{b}}$، والترتيب فيه
\emph{معتبر}؛ فإذا كان $\OLId{latin-ordinary}{a} \neq \OLId{latin-ordinary}{b}$ كان $\tuple{\OLId{latin-ordinary}{a}, \OLId{latin-ordinary}{b}} \neq \tuple{\OLId{latin-ordinary}{b}, \OLId{latin-ordinary}{a}}$.
وهذا بخلاف الزوج غير المرتب، أي المجموعة التي عدد عناصرها $\OLMathNumeral{2}$، إذ
$\{\OLId{latin-ordinary}{a}, \OLId{latin-ordinary}{b}\}=\{\OLId{latin-ordinary}{b}, \OLId{latin-ordinary}{a}\}$. والآخر \emph{الضرب الديكارتي}: فالمجموعتان
$\OLId{latin-looped}{A}$ و$\OLId{latin-looped}{B}$ تعطيان~$\OLId{latin-looped}{A} \times \OLId{latin-looped}{B}$، وهي مجموعة كل زوج $\tuple{\OLId{latin-ordinary}{x}, \OLId{latin-ordinary}{y}}$
حيث $\OLId{latin-ordinary}{x} \in \OLId{latin-looped}{A}$ و$\OLId{latin-ordinary}{y} \in \OLId{latin-looped}{B}$. وبخاصة، $\OLId{latin-looped}{A}^{\OLMathNumeral{2}}= \OLId{latin-looped}{A} \times \OLId{latin-looped}{A}$ يجمع
الأزواج المرتبة من~$\OLId{latin-looped}{A}$ كلها.

ولنخص بالنظر العلاقة $<$ على مجموعة الأعداد الطبيعية~$\Nat$.
اجمع أزواج الأعداد $\tuple{\OLId{latin-ordinary}{n}, \OLId{latin-ordinary}{m}}$ التي تحقق $\OLId{latin-ordinary}{n}<\OLId{latin-ordinary}{m}$ في المجموعة:
\[
  \OLId{latin-looped}{R}=\Setabs{\tuple{\OLId{latin-ordinary}{n}, \OLId{latin-ordinary}{m}}}{\OLId{latin-ordinary}{n}, \OLId{latin-ordinary}{m} \in \Nat \text{ و } \OLId{latin-ordinary}{n}<\OLId{latin-ordinary}{m}}.
\]
فكون $\OLId{latin-ordinary}{n}$ أصغر من $\OLId{latin-ordinary}{m}$ وعضوية الزوج $\tuple{\OLId{latin-ordinary}{n}, \OLId{latin-ordinary}{m}}$ في~$\OLId{latin-looped}{R}$
متكافئان، أي:
\[
      \OLId{latin-ordinary}{n}<\OLId{latin-ordinary}{m}\text{ إذا وفقط إذا }\tuple{\OLId{latin-ordinary}{n}, \OLId{latin-ordinary}{m}} \in \OLId{latin-looped}{R}.
\]
فلا يفوتنا شيء من المعلومات إن جعلنا $\OLId{latin-looped}{R}$ \emph{نفس} العلاقة $<$
على~$\Nat$.

وعلى هذا الوجه نمثل كل علاقة بين الأعداد بجزئية من $\Nat^{\OLMathNumeral{2}}$.
والعكس كذلك: فكل مجموعة أزواج $\OLId{latin-looped}{S} \subseteq \Nat^{\OLMathNumeral{2}}$ تقابلها علاقة
بين الأعداد، تقوم بين $\OLId{latin-ordinary}{n}$ و$\OLId{latin-ordinary}{m}$ إذا وفقط إذا كان $\tuple{\OLId{latin-ordinary}{n}, \OLId{latin-ordinary}{m}} \in \OLId{latin-looped}{S}$.
ومن هذا يصح التعريف الآتي.
\end{explain}

\begin{defn}[العلاقة الثنائية]
\emph{العلاقة الثنائية} على المجموعة $\OLId{latin-looped}{A}$ كل مجموعة جزئية من~$\OLId{latin-looped}{A}^{\OLMathNumeral{2}}$.
فإذا كانت $\OLId{latin-looped}{R} \subseteq \OLId{latin-looped}{A}^{\OLMathNumeral{2}}$ علاقة على~$\OLId{latin-looped}{A}$، و$\OLId{latin-ordinary}{x}, \OLId{latin-ordinary}{y} \in \OLId{latin-looped}{A}$،
جاز أن نكتب $Rxy$ (أو $xRy$) اختصارًا لـ$\tuple{\OLId{latin-ordinary}{x}, \OLId{latin-ordinary}{y}} \in \OLId{latin-looped}{R}$.
\end{defn}

\begin{ex}
  \ollabel{relations}
لنرتب أزواج الأعداد الطبيعية $\Nat^{\OLMathNumeral{2}}$ في مصفوفة ذات بعدين:
\[
  \begin{array}{ccccc}
  \mathbf{\tuple{ 0,0 }} & \tuple{ \OLMathNumeral{0},\OLMathNumeral{1} } &
    \tuple{ \OLMathNumeral{0},\OLMathNumeral{2} } & \tuple{ \OLMathNumeral{0},\OLMathNumeral{3} } & \ldots\\
  \tuple{ \OLMathNumeral{1},\OLMathNumeral{0} } & \mathbf{\tuple{ 1,1 }} &
    \tuple{ \OLMathNumeral{1},\OLMathNumeral{2} } & \tuple{ \OLMathNumeral{1},\OLMathNumeral{3} } & \ldots\\
  \tuple{ \OLMathNumeral{2},\OLMathNumeral{0} } & \tuple{ \OLMathNumeral{2},\OLMathNumeral{1} } &
    \mathbf{\tuple{ 2,2 }} & \tuple{ \OLMathNumeral{2},\OLMathNumeral{3} } & \ldots\\
  \tuple{ \OLMathNumeral{3},\OLMathNumeral{0} } & \tuple{ \OLMathNumeral{3},\OLMathNumeral{1} } & \tuple{ \OLMathNumeral{3},\OLMathNumeral{2} } &
    \mathbf{\tuple{ 3,3 }} & \ldots\\
  \vdots & \vdots & \vdots & \vdots & \mathbf{\ddots}
  \end{array}
\]
وقد ميّز خط القطر بالعريض؛ فإن جزئية $\Nat^\OLMathNumeral{2}$ التي تضم أزواجه، أي:
\[
  \{\tuple{\OLMathNumeral{0},\OLMathNumeral{0} }, \tuple{ \OLMathNumeral{1},\OLMathNumeral{1} }, \tuple{ \OLMathNumeral{2},\OLMathNumeral{2} }, \dots\},
  \]
هي \emph{علاقة الهوية على}~$\Nat$. ولشيوع هذه العلاقة نجعل
$\Id{\OLId{latin-looped}{A}}=\Setabs{\tuple{ \OLId{latin-ordinary}{x},\OLId{latin-ordinary}{x} }}{\OLId{latin-ordinary}{x} \in \OLId{latin-looped}{A}}$ لكل مجموعة $\OLId{latin-looped}{A}$.
وأما الأزواج التي فوق القطر، فمجموعتها:
\[
  \OLId{latin-looped}{L} = \{\tuple{ \OLMathNumeral{0},\OLMathNumeral{1} },\tuple{ \OLMathNumeral{0},\OLMathNumeral{2} },\ldots,\tuple{ \OLMathNumeral{1},\OLMathNumeral{2} },
  \tuple{ \OLMathNumeral{1},\OLMathNumeral{3} }, \dots, \tuple{ \OLMathNumeral{2},\OLMathNumeral{3} }, \tuple{ \OLMathNumeral{2},\OLMathNumeral{4} },\ldots\},
\]
وهي علاقة \emph{أصغر من}؛ فـ$Lnm$ يكافئ $\OLId{latin-ordinary}{n}<\OLId{latin-ordinary}{m}$.
وأما ما تحت القطر فمجموعته:
\[
  \OLId{latin-looped}{G}=\{ \tuple{ \OLMathNumeral{1},\OLMathNumeral{0} },\tuple{ \OLMathNumeral{2},\OLMathNumeral{0} },\tuple{
    \OLMathNumeral{2},\OLMathNumeral{1} }, \tuple{ \OLMathNumeral{3},\OLMathNumeral{0} },\tuple{ \OLMathNumeral{3},\OLMathNumeral{1} },\tuple{ \OLMathNumeral{3},\OLMathNumeral{2} }, \dots\},
\]
وهي علاقة \emph{أكبر من}؛ فـ$Gnm$ يكافئ $\OLId{latin-ordinary}{n}>\OLId{latin-ordinary}{m}$. وإذا ضممنا
$\OLId{latin-looped}{L}$ إلى $\OLId{latin-looped}{I}$، حيث $\OLId{latin-looped}{I} = \Id{\Nat}$، وسمينا الاتحاد $\OLId{latin-looped}{K}=\OLId{latin-looped}{L}\cup \OLId{latin-looped}{I}$ حصلت علاقة \emph{أصغر من أو
يساوي}، إذ $Knm$ يكافئ $\OLId{latin-ordinary}{n} \le \OLId{latin-ordinary}{m}$. وكذلك $\OLId{latin-looped}{H}=\OLId{latin-looped}{G} \cup \OLId{latin-looped}{I}$ هي
\emph{علاقة أكبر من أو يساوي.} فهذه العلاقات $\OLId{latin-looped}{L}$ و$\OLId{latin-looped}{G}$ و$\OLId{latin-looped}{K}$ و$\OLId{latin-looped}{H}$
صنف مخصوص يسمى \emph{الترتيبات}. وأما $\OLId{latin-looped}{L}$ و$\OLId{latin-looped}{G}$ فلا تقوم واحدة منهما
بين عدد ونفسه: فلا $\OLId{latin-looped}{L}$ ولا $\OLId{latin-looped}{G}$ تربط العدد بنفسه؛ أي لكل $\OLId{latin-ordinary}{n}$ يمتنع
$Lnn$ و$Gnn$. وهذه الخاصية تسمى \emph{اللاانعكاسية}، فإذا اتصف
بها ترتيب سمي \emph{ترتيبًا صارمًا.}
\end{ex}

\textbf{تنبيه تصحيحي على الأصل (C012-01).} صُرّح هنا بأن $\OLId{latin-looped}{I}$ علاقة الهوية على
$\Nat$؛ فقد استعملها الأصل بهذا المعنى من غير أن يثبت لها هذا الاختصار.

\begin{explain}
ولئن كانت الترتيبات والهوية من العلاقات الطبيعية المهمة، فتعريفنا
أعم منها؛ إذ \emph{كل} جزئية من $\OLId{latin-looped}{A}^{\OLMathNumeral{2}}$ علاقة على~$\OLId{latin-looped}{A}$، وإن بدت
مصطنعة بعيدة عن المألوف. فـ$\emptyset$ علاقة على كل مجموعة، تسمى
\emph{العلاقة الخالية}، ولا تقوم بين زوج من العناصر. و$\OLId{latin-looped}{A}^{\OLMathNumeral{2}}$ نفسها
علاقة على~$\OLId{latin-looped}{A}$، تقوم بين كل زوج، واسمها \emph{العلاقة الشاملة}.
بل تعد المجموعة
$\OLId{latin-looped}{E}=\Setabs{\tuple{\OLId{latin-ordinary}{n}, \OLId{latin-ordinary}{m}}}{\OLId{latin-ordinary}{n}>\OLMathNumeral{5} \text{ أو } \OLId{latin-ordinary}{m} \times \OLId{latin-ordinary}{n} \ge \OLMathNumeral{34}}$ علاقة أيضًا.
\end{explain}

\begin{prob}
عيّن !!{element}s العلاقة $\subseteq$ على مجموعة $\Pow{\{\OLId{latin-ordinary}{a}, \OLId{latin-ordinary}{b}, \OLId{latin-ordinary}{c}\}}$.
\end{prob}

\end{document}
